

\begin{theorem}
Let
\[
s=\frac12+i\gamma,\qquad a=1-s=\frac12-i\gamma,
\]
and assume
\[
\zeta(s)=0,\qquad \gamma>\frac{\pi}{\log 2}.
\]
For \(N\in\mathbb Z_{\ge 0}\), define
\[
S_\gamma(N):=\sum_{n=1}^N n^{-s},\qquad S_\gamma(0):=0.
\]
For \(p,t\in\mathbb Z_{\ge 1}\), set
\[
D_\gamma(t,p):=S_\gamma(t)-p^a\,S_\gamma\!\left(\Bigl\lfloor \frac{t}{p}\Bigr\rfloor\right),
\]
and
\[
\tau_p(\gamma):=\inf\{t\ge 1:\ |D_\gamma(t,p)|\ge \sqrt p\}\in \mathbb Z_{\ge 1}\cup\{\infty\}.
\]

Then there exists a unique \(\phi_\gamma\in(0,\pi)\) such that
\[
e^{\phi_\gamma/(2\gamma)}=\cos\phi_\gamma+2\gamma\sin\phi_\gamma.
\]
If
\[
\kappa_\gamma:=e^{\phi_\gamma/\gamma},
\]
then
\[
\lim_{p\to\infty}\frac{\tau_p(\gamma)}{p}=\kappa_\gamma.
\]

Moreover, as \(\gamma\to\infty\),
\[
\phi_\gamma=\pi-\frac1\gamma+O(\gamma^{-2}),
\qquad
\gamma\log \kappa_\gamma=\pi-\frac1\gamma+O(\gamma^{-2}).
\]
Consequently, for every sequence \((\gamma_n)\) such that
\[
\zeta\!\left(\frac12+i\gamma_n\right)=0,\qquad \gamma_n\to\infty,
\]
one has
\[
\lim_{n\to\infty}\gamma_n\log\!\left(\lim_{p\to\infty}\frac{\tau_p(\gamma_n)}{p}\right)=\pi.
\]
No simplicity hypothesis on the zero is used.
\end{theorem}

\begin{proof}
For \(N\ge 1\), Euler--Maclaurin gives
\[
\sum_{n=1}^N n^{-s}
=
\zeta(s)+\frac{N^{1-s}}{1-s}+\frac12 N^{-s}+O_\gamma(N^{-3/2}).
\]
Since \(\zeta(s)=0\) and \(a=1-s\),
\begin{equation}
S_\gamma(N)=\frac{N^a}{a}+\frac12 N^{-s}+O_\gamma(N^{-3/2}).
\label{eq:EM}
\end{equation}

If \(t<p\), then \(\lfloor t/p\rfloor=0\), hence \(D_\gamma(t,p)=S_\gamma(t)\). By
\eqref{eq:EM},
\[
|D_\gamma(t,p)|
\le \frac{\sqrt t}{|a|}+C_\gamma
\le \frac{\sqrt p}{|a|}+C_\gamma<\sqrt p
\]
for all sufficiently large \(p\), because \(|a|=\sqrt{\frac14+\gamma^2}>1\). Thus
\begin{equation}
\tau_p(\gamma)\ge p
\qquad (p\to\infty).
\label{eq:before-p}
\end{equation}

Now suppose \(p\le t<2p\). Then \(\lfloor t/p\rfloor=1\), so
\[
D_\gamma(t,p)=S_\gamma(t)-p^a.
\]
Write \(\kappa=t/p\in[1,2)\). Using \eqref{eq:EM},
\[
p^{-a}D_\gamma(t,p)
=
\frac{\kappa^a}{a}-1+\frac{1}{2p}\kappa^{-s}+O_\gamma(p^{-2}),
\]
uniformly for \(\kappa\in[1,2)\). Therefore
\begin{equation}
\left|
\frac{|D_\gamma(t,p)|}{\sqrt p}-F_\gamma(\kappa)
\right|
\le \frac{C_\gamma}{p},
\qquad
F_\gamma(\kappa):=\left|\frac{\kappa^a}{a}-1\right|.
\label{eq:profile}
\end{equation}

Set
\[
A:=|a|^2=\frac14+\gamma^2,
\qquad
\phi:=\gamma\log\kappa.
\]
Then \(\kappa^a=e^{\phi/(2\gamma)}e^{-i\phi}\), and a direct calculation yields
\begin{equation}
F_\gamma(\kappa)^2-1
=
\frac{e^{\phi/(2\gamma)}}{A}
\Bigl(e^{\phi/(2\gamma)}-\cos\phi-2\gamma\sin\phi\Bigr).
\label{eq:F-formula}
\end{equation}
Define
\[
Q_\gamma(\phi):=e^{-\phi/(2\gamma)}(\cos\phi+2\gamma\sin\phi).
\]
By \eqref{eq:F-formula},
\[
F_\gamma(\kappa)<1 \iff Q_\gamma(\phi)>1,
\qquad
F_\gamma(\kappa)=1 \iff Q_\gamma(\phi)=1.
\]

Moreover,
\[
Q_\gamma'(\phi)
=
e^{-\phi/(2\gamma)}
\left(
\left(2\gamma-\frac{1}{2\gamma}\right)\cos\phi-2\sin\phi
\right).
\]
Since \(2\gamma-\frac1{2\gamma}>0\), the bracket is of the form
\(R\cos(\phi+\alpha)\) with \(R>0\) and \(\alpha\in(0,\pi/2)\). Hence
\(Q_\gamma'\) has exactly one zero in \((0,\pi)\). As
\[
Q_\gamma(0)=1,\qquad
Q_\gamma'(0)=2\gamma-\frac1{2\gamma}>0,\qquad
Q_\gamma(\pi)=-e^{-\pi/(2\gamma)}<1,
\]
there is a unique \(\phi_\gamma\in(0,\pi)\) such that
\[
Q_\gamma(\phi_\gamma)=1,
\]
equivalently
\begin{equation}
e^{\phi_\gamma/(2\gamma)}
=
\cos\phi_\gamma+2\gamma\sin\phi_\gamma.
\label{eq:def-phi}
\end{equation}
Set
\[
\kappa_\gamma:=e^{\phi_\gamma/\gamma}.
\]
Since \(\phi_\gamma<\pi\) and \(\gamma>\pi/\log 2\),
\[
1<\kappa_\gamma<e^{\pi/\gamma}<2.
\]
Also,
\begin{equation}
F_\gamma(\kappa)<1 \quad (1<\kappa<\kappa_\gamma),
\qquad
F_\gamma(\kappa)>1 \quad (\kappa_\gamma<\kappa<e^{\pi/\gamma}).
\label{eq:sign}
\end{equation}

We now handle the discrete neighborhood of \(\kappa=1\). Let
\[
z_0:=\frac1a-1=\frac{\frac12+i\gamma}{\frac12-i\gamma}.
\]
Then \(|z_0|=1\) and
\[
\Re z_0=\frac{\frac14-\gamma^2}{\frac14+\gamma^2}<0.
\]
Fix \(h\in\mathbb Z_{\ge 0}\), and put \(t=p+h\). From \eqref{eq:EM},
\[
p^{-a}D_\gamma(p+h,p)
=
\frac{(1+h/p)^a}{a}-1+\frac{1}{2p}(1+h/p)^{-s}+O_{\gamma,h}(p^{-2}),
\]
so
\begin{equation}
p^{-a}D_\gamma(p+h,p)
=
z_0+\frac{h+\frac12}{p}+O_{\gamma,h}(p^{-2}).
\label{eq:local}
\end{equation}
Therefore
\[
\frac{|D_\gamma(p+h,p)|^2}{p}
=
1+\frac{2\Re z_0\,(h+\frac12)}{p}+O_{\gamma,h}(p^{-2})<1
\]
for all sufficiently large \(p\). Thus, for every fixed \(h\ge 0\),
\begin{equation}
|D_\gamma(p+h,p)|<\sqrt p
\qquad (p\to\infty).
\label{eq:fixed-h}
\end{equation}

Next, \(F_\gamma(1)=1\), and differentiating \(F_\gamma(\kappa)=|\kappa^a/a-1|\) at
\(\kappa=1\) gives
\[
F_\gamma'(1)=\Re z_0<0.
\]
Hence there exist \(\delta_\gamma,c_\gamma>0\) such that
\begin{equation}
F_\gamma(1+u)\le 1-c_\gamma u
\qquad (0\le u\le \delta_\gamma).
\label{eq:linear-gap}
\end{equation}

Fix \(\varepsilon>0\) so small that
\[
0<\varepsilon<\frac12\min\{\kappa_\gamma-1,\ \delta_\gamma,\ e^{\pi/\gamma}-\kappa_\gamma\}.
\]
Choose \(M\in\mathbb N\) so large that \(M>2C_\gamma/c_\gamma\). If
\(M\le h\le \varepsilon p\), then \(u=h/p\in[0,\delta_\gamma]\), and
\eqref{eq:profile} together with \eqref{eq:linear-gap} gives
\[
\frac{|D_\gamma(p+h,p)|}{\sqrt p}
\le
F_\gamma(1+h/p)+\frac{C_\gamma}{p}
\le
1-c_\gamma\frac{h}{p}+\frac{C_\gamma}{p}
<1.
\]
Combining this with \eqref{eq:fixed-h}, we obtain
\begin{equation}
|D_\gamma(t,p)|<\sqrt p
\qquad
\bigl(p\le t\le (1+\varepsilon)p\bigr),
\label{eq:first-window}
\end{equation}
for all sufficiently large \(p\).

Since \(F_\gamma<1\) on \((1,\kappa_\gamma)\), compactness and \eqref{eq:sign} imply that
there exists \(\eta_\gamma(\varepsilon)>0\) such that
\[
F_\gamma(\kappa)\le 1-\eta_\gamma(\varepsilon)
\qquad
(1+\varepsilon\le \kappa\le \kappa_\gamma-\varepsilon).
\]
Using \eqref{eq:profile} once again,
\begin{equation}
|D_\gamma(t,p)|<\sqrt p
\qquad
\bigl((1+\varepsilon)p\le t\le (\kappa_\gamma-\varepsilon)p\bigr),
\label{eq:middle-window}
\end{equation}
for all sufficiently large \(p\). Together with \eqref{eq:before-p},
\eqref{eq:first-window}, and \eqref{eq:middle-window}, this yields
\begin{equation}
\tau_p(\gamma)>(\kappa_\gamma-\varepsilon)p
\qquad (p\to\infty).
\label{eq:lower}
\end{equation}

For the upper bound, the interval
\[
\left[\kappa_\gamma+\frac{\varepsilon}{2},\,\kappa_\gamma+\varepsilon\right]
\]
is compactly contained in \((\kappa_\gamma,e^{\pi/\gamma})\). Hence, by
\eqref{eq:sign}, there exists \(\eta_\gamma'(\varepsilon)>0\) such that
\[
F_\gamma(\kappa)\ge 1+\eta_\gamma'(\varepsilon)
\qquad
\left(\kappa_\gamma+\frac{\varepsilon}{2}\le \kappa\le \kappa_\gamma+\varepsilon\right).
\]
Let
\[
t_+(p):=\left\lceil\left(\kappa_\gamma+\frac{\varepsilon}{2}\right)p\right\rceil.
\]
For all sufficiently large \(p\),
\[
\frac{t_+(p)}{p}\in
\left[\kappa_\gamma+\frac{\varepsilon}{2},\,\kappa_\gamma+\varepsilon\right],
\]
so \eqref{eq:profile} gives
\[
|D_\gamma(t_+(p),p)|>\sqrt p.
\]
Therefore
\begin{equation}
\tau_p(\gamma)\le t_+(p)\le (\kappa_\gamma+\varepsilon)p
\qquad (p\to\infty).
\label{eq:upper}
\end{equation}

Since \(\varepsilon>0\) is arbitrary, \eqref{eq:lower} and \eqref{eq:upper} imply
\[
\lim_{p\to\infty}\frac{\tau_p(\gamma)}{p}=\kappa_\gamma.
\]

It remains to analyze \(\kappa_\gamma\) for large \(\gamma\). From
\eqref{eq:def-phi},
\[
Q_\gamma(\pi-\delta)
=
e^{-(\pi-\delta)/(2\gamma)}
\bigl(-\cos\delta+2\gamma\sin\delta\bigr)\to+\infty
\qquad (\gamma\to\infty)
\]
for each fixed \(\delta\in(0,\pi)\). Hence \(\phi_\gamma\to\pi\). Write
\[
\varepsilon_\gamma:=\pi-\phi_\gamma.
\]
Then \(\varepsilon_\gamma\to 0\), and \eqref{eq:def-phi} becomes
\begin{equation}
2\gamma\sin\varepsilon_\gamma
=
e^{(\pi-\varepsilon_\gamma)/(2\gamma)}+\cos\varepsilon_\gamma.
\label{eq:eps}
\end{equation}
The right-hand side is \(O(1)\), so \(\gamma\varepsilon_\gamma=O(1)\), hence
\(\varepsilon_\gamma=O(\gamma^{-1})\). Expanding \eqref{eq:eps},
\[
2\gamma\varepsilon_\gamma+O(\gamma\varepsilon_\gamma^3)
=
1+\frac{\pi-\varepsilon_\gamma}{2\gamma}+O(\gamma^{-2})
+
1-\frac{\varepsilon_\gamma^2}{2}+O(\varepsilon_\gamma^4).
\]
Because \(\varepsilon_\gamma=O(\gamma^{-1})\), all discarded terms are
\(O(\gamma^{-2})\), so
\[
2\gamma\varepsilon_\gamma
=
2+\frac{\pi}{2\gamma}+O(\gamma^{-2}).
\]
In particular,
\[
\varepsilon_\gamma=\frac1\gamma+O(\gamma^{-2}),
\]
and therefore
\[
\phi_\gamma=\pi-\frac1\gamma+O(\gamma^{-2}).
\]
Since \(\kappa_\gamma=e^{\phi_\gamma/\gamma}\),
\[
\gamma\log\kappa_\gamma=\phi_\gamma
=
\pi-\frac1\gamma+O(\gamma^{-2}).
\]

The final assertion follows immediately by applying this estimate to any zero
sequence \(\gamma_n\to\infty\).
\end{proof}
